%% kapa: exam=1 round=31 year=2020 period=2 form=A subject=accounting subject_ko=회계학 seq=14 number=54
%% source: Q-Net 감정평가사 기출문제 (cst00309) artlSeq=5207219 file=extracted/1008/5207219/A형/A형/2020년도 제31회 감정평가사 1차 시험_2교시_A형.pdf
%% conversion: kordoc 4.24.51 -> tools/gen1.py (auto); review=proofread-sonnet; flags=none
\begin{question}{54}
(주)감평은 20x1년 1월 1일 제품을 판매하기로 (주)한국과 계약을 체결하였다. 동 제품에 대한 통제는 20x2년 말에 (주)한국으로 이전된다. 계약에 의하면 (주)한국은 ㉠\uline{계약을 체결할 때 ￦100,000을 지급}하거나 ㉡\uline{제품을 통제하는 20x2년 말에 ￦125,440을 지급}하는 방법 중 하나를 선택할 수 있다. 이 중 (주)한국은 ㉠을 선택함으로써 계약체결일에 현금 ￦100,000을 (주)감평에게 지급하였다. (주)감평은 자산 이전시점과 고객의 지급시점 사이의 기간을 고려하여 유의적인 금융요소가 포함되어 있다고 판단하고 있으며, (주)한국과 별도 금융거래를 한다면 사용하게 될 증분차입이자율 연 10\%를 적절한 할인율로 판단한다. 동 거래와 관련하여 (주)감평이 20x1년 말 재무상태표에 계상할 계약부채의 장부금액(A)과 20x2년도 포괄손익계산서에 인식할 매출수익(B)은?
{\par\noindent\hspace*{2em}\makebox[3cm][c]{\uline{(A)}}\makebox[3cm][c]{\uline{(B)}}\par}
\begin{choices}
  \item \makebox[3cm][c]{￦100,000}\makebox[3cm][c]{￦100,000}
  \item \makebox[3cm][c]{￦110,000}\makebox[3cm][c]{￦121,000}
  \item \makebox[3cm][c]{￦110,000}\makebox[3cm][c]{￦125,440}
  \item \makebox[3cm][c]{￦112,000}\makebox[3cm][c]{￦121,000}
  \item \makebox[3cm][c]{￦112,000}\makebox[3cm][c]{￦125,440}
\end{choices}
\end{question}
\answer{2}
